\section{Open problems}\label{sec:open}

Each problem is stated with the reason it is open, so that a reader can tell a hard problem from
an unattempted one.

\begin{enumerate}[label=\textbf{O\arabic*.},leftmargin=2.4em]
\item \textbf{The angle over a long path.} Theorem~\ref{thm:riesz} settles the derivative bound
(constant $1$, sharp), gives the union-gap form for a single perturbation, and identifies $\pi/2$
as the endpoint value of the angle bound. What is not settled is the \emph{integrated} version:
along a long path with a gap that may shrink, the angle can pass through $\pi/2$ and the count can
change, and no bound is known for the total angular travel in terms of
$\int\norm{\dot A}/\Gap\,\df t$. This is the adiabatic question in a finite-dimensional setting,
and it is open because the derivative bound cannot be integrated past the point where the angle
reaches $\pi/2$. The sketch of what a proof would need is in
\gate{notes/04\_PI\_OVER\_2.md}, whose earlier version stated the problem in the form that
Theorem~\ref{thm:riesz} has now resolved.

\item \textbf{The rank-$k$ $\tan 2\Theta$ bound.} Theorem~\ref{thm:tan2} settles rank one
completely, with an exact constant and an exact sharpness family. The gap to close is
$k\ge2$: given a cluster projector $\PP$ of rank $k$ and an off-diagonal perturbation $V$ whose
spectral sets have disjoint convex hulls, is $\abs{\tan\Theta}\le\delta/d$ for the operator
angle, and with which constant? The rank-one proof does not extend, and the obstruction is
identifiable: its decisive step is $b-a\ge d$, where $b=\ip{Aw}{w}$ ranges over the numerical
range of the compression; for $k\ge2$ one gets such a two-dimensional reduction for each
eigenvector separately, but they do not assemble into a bound on the angular operator
$X$ because the Riccati equation is quadratic, and the elementary trace estimate gives only
$w\norm{X}_F^2-d\norm{X}_F+w\ge0$, which is vacuous when $w<d/2$. What is needed is a solution
theory for $A_+X-XA_-=XBX-B^*$ with an a priori bound on $\norm X$: this is precisely the
machinery of \cite{kmm1,kmm2,kmm3}, which we cite and have not re-derived, and which is where a
proof would begin. A gate sampling the rank-$k$ separated case is described in
\gate{notes/06\_CONVEX\_HULL.md}.

\item \textbf{The optic points, exactly.} Theorem~\ref{thm:optic} locates the two optic points
numerically and gives the displacement law $\varepsilon^2/9+O(\varepsilon^3)$ for one family.
What is missing is a closed form: the displacement is the zero of a two-by-two system whose
entries are algebraic in the masses, so a formula exists and has not been found. The same
question for $N\ge4$ bodies is open and is the natural continuation of \S\ref{sec:shape}.

\item \textbf{A per-shape capacity law.} The sum rule fixes the total; it says nothing about how
the total is distributed. A distribution law would need a statement about the eigenvalues of
$M^*M$ beyond their sum, i.e.\ beyond Theorems~\ref{thm:shape-sum}--\ref{thm:shape-det}. For the
three-body problem the pair $(\operatorname{tr},\det)$ already determines the spectrum, so the
problem is closed there and open from $N=4$ on.

\item \textbf{The square case in general dimension.} When $N=d+1$ the Gram matrix $G=X^\top X$ is
square, its trace-one slice is the shape space, and the transport operator splits into an
eigenvalue sector and an eigenframe sector whose spectra are computable. The corpus states this
with an interlacing theorem; the present paper does not use it, and the statement has not been
re-derived here. It is the natural next theorem, and it is the one most in need of being written
down with its hypotheses explicit.

\item \textbf{The census, quantitatively.} The corpus counts channels and measures plateaus with
a fitted constant $C_0$. Given Theorem~\ref{thm:sumrule}, the natural question is exact: for
$N$ bodies, how many pairwise distances can be resolved to accuracy $\varepsilon$ as the shape
varies over a ball of radius $\delta$? That is a question about the singular values of $M$ and
hence about the constant in a $\sin\theta\ge\varepsilon$ dictionary; it is open, and it is the
form in which ``the price of a direction'' can become a theorem rather than a measurement.
\end{enumerate}
